\documentclass[11pt]{article}
\usepackage{certmachine}
%% generated by tools/paper-numbers/navier-stokes-audit.js from corpus/navier-stokes/{lean-repo,build,probes,audit,nonvacuity,MANIFEST}.json, corpus/navier-stokes/euler.md and certs/claims-ledger.json — do not edit (git eece01c1497a)
\newcommand{\RepoCommit}{eece01c1497a}
\newcommand{\LeanCommit}{8937a8f4}
\newcommand{\LeanCommitFull}{8937a8f4cbc7abaab5e9e97d1cc7f5d2319d9538}
\newcommand{\LeanCommitDate}{2026-09-08}
\newcommand{\LeanCommitTime}{10:57 UTC}
\newcommand{\LeanToolchain}{v4.34.0-rc2}
\newcommand{\MathlibRev}{85e3a25e}
\newcommand{\MathlibRevFull}{85e3a25e006c35636f0e53b0e9296caca2685bc0}
\newcommand{\ComparatorRev}{19e111e2}
\newcommand{\LeanExportRev}{cacf989b}
\newcommand{\FormalConjRev}{8bf45ed7}
\newcommand{\FormalConjTouched}{2026-07-27}
\newcommand{\FormalConjAdded}{2026-05-15}
\newcommand{\FormalConjPR}{1457}
\newcommand{\CorpusFiles}{24}
\newcommand{\CorpusFetched}{2026-09-09}
\newcommand{\PaperSha}{0e779481}
\newcommand{\PaperPages}{166}
\newcommand{\EulerPages}{57}
\newcommand{\LeanFiles}{2{,}486}
\newcommand{\LeanLines}{618{,}762}
\newcommand{\NSFiles}{644}
\newcommand{\NSLines}{404{,}869}
\newcommand{\EulerFiles}{1{,}840}
\newcommand{\EulerLines}{213{,}419}
\newcommand{\ChalFiles}{2}
\newcommand{\ChalLines}{474}
\newcommand{\HygSorryChal}{4}
\newcommand{\HygPrintAxioms}{4}
\newcommand{\HygHeartbeatComments}{5}
\newcommand{\UpstreamDiffLines}{4}
\newcommand{\FormalReview}{self-assessed}
\newcommand{\FormalAutomation}{GPT-6 Astra}
\newcommand{\BuildMachine}{Mac14,2}
\newcommand{\BuildCpus}{8}
\newcommand{\BuildMemory}{16}
\newcommand{\BuildOS}{darwin 25.4.0}
\newcommand{\BuildArch}{arm64}
\newcommand{\BuildDate}{2026-09-09}
\newcommand{\BuildMinutes}{140}
\newcommand{\BuildJobs}{2{,}486}
\newcommand{\BuildJobsTotal}{11{,}251}
\newcommand{\BuildLegOne}{681}
\newcommand{\BuildLegOneThreads}{8}
\newcommand{\BuildLegTwo}{1{,}805}
\newcommand{\BuildLegTwoThreads}{4}
\newcommand{\BuildExit}{0}
\newcommand{\SlowestModule}{Euler.HeatRegularizedPaths}
\newcommand{\SlowestMinutes}{14}
\newcommand{\AxiomDecls}{4}
\newcommand{\AxiomCount}{3}
\newcommand{\AxiomList}{\texttt{propext}, \texttt{Classical.choice}, \texttt{Quot.sound}}
\newcommand{\ComparatorConfigs}{2}
\newcommand{\ComparatorAccepted}{2}
\newcommand{\ComparatorKernelAccepts}{4}
\newcommand{\ComparatorMinutesNS}{13}
\newcommand{\ComparatorMinutesEuler}{17}
\newcommand{\NonVacuityDecls}{7}
\newcommand{\NonVacuityErrors}{0}
\newcommand{\NonVacuityFile}{cert-machine-NonVacuity.lean}
\newcommand{\FidelityRows}{9}
\newcommand{\FidelitySame}{7}
\newcommand{\FidelityEquivalent}{2}
\newcommand{\PaperVsLeanRows}{7}
\newcommand{\PaperClausesNotFormal}{2}
\newcommand{\Readers}{six}
\newcommand{\SectionMemos}{five}
\newcommand{\VerifiedRows}{6}
\newcommand{\ConeLoopTrials}{371}
\newcommand{\ConeBreakPct}{7}
\newcommand{\WeakPoints}{13}
\newcommand{\WeakPointsAfterZero}{12}
\newcommand{\Slips}{5}
\newcommand{\TracedFull}{6}
\newcommand{\TracedStronger}{2}
\newcommand{\TracedAltered}{5}
\newcommand{\TracedNone}{1}
\newcommand{\SlipsProseOnly}{3}
\newcommand{\ObstructionRows}{9}
\newcommand{\ObstructionsMet}{6}
\newcommand{\ObstructionsEvaded}{3}
\newcommand{\ProbeScripts}{7}
\newcommand{\ProbeReds}{26}
\newcommand{\ProbeRedsDead}{0}
\newcommand{\CoeffIdentities}{69}
\newcommand{\ConeSweepPoints}{25{,}200}
\newcommand{\ConeAdmissible}{7{,}234}
\newcommand{\ConeNotAdmissible}{17{,}966}
\newcommand{\ConeSignFlipDisagree}{10{,}456}
\newcommand{\ConeRefPoints}{1{,}640}
\newcommand{\ConeRefInside}{224}
\newcommand{\HeatResidual}{$4.7\times10^{-19}$}
\newcommand{\HeatTaylorOrder}{4}
\newcommand{\AxisLambdaFrom}{10^{6}}
\newcommand{\HeatDigits}{25}
\newcommand{\EulerChecksUnrun}{3}
\newcommand{\LedgerId}{navier-stokes-openai-2026}
\newcommand{\LedgerVerdict}{PARTIAL}
\newcommand{\LedgerRecorded}{2026-09-09}
\newcommand{\DateRowsCount}{11}
\newcommand{\DateRows}{2026-05-15 & DeepMind's Formal Conjectures adds the Lean statement of the four Clay alternatives (\#1457) \\
2026-07-27 & last change to that file (\#4643, divergence lemmas); the revision OpenAI pinned \\
2026-08-15 & Alp\"oge--Buckmaster obtain smooth-forcing blowup for Boussinesq and Euler (their statement) \\
2026-08-22 & Alp\"oge--Buckmaster's Lean verification (their statement) \\
2026-08-28 & OpenAI begins training the internal model used (the announcement) \\
2026-09-01 & OpenAI launches the agent effort after 'a rumor that two Millennium Prize problems had been resolved' (the announcement) \\
2026-09-03 & Tao's six-part post on the opportunity cost of an autonomous solution; Buckmaster writes to OpenAI (his statement) \\
2026-09-05 & OpenAI's agents 'arrive at their resolution', 88 hours in (the announcement); Tao's clarification: 'not aware of any significant developments' \\
2026-09-06 & Lean formalization complete, 17 hours later (the announcement); the two calls with Buckmaster (his statement) \\
2026-09-08 03:58 UTC & Alp\"oge--Buckmaster publish IPM, Boussinesq and Euler with smooth forcing, Lean at tristanbuckmaster/fluid\_lean (repository created 04:03 UTC) \\
2026-09-08 10:53 UTC & openai/NavierStokesAndEuler created; its single commit at 10:57 UTC; the announcement the same day \\}
\newcommand{\DateList}{\textbf{2026-05-15}: DeepMind's Formal Conjectures adds the Lean statement of the four Clay alternatives (\#1457); \textbf{2026-07-27}: last change to that file (\#4643, divergence lemmas); the revision OpenAI pinned; \textbf{2026-08-15}: Alp\"oge--Buckmaster obtain smooth-forcing blowup for Boussinesq and Euler (their statement); \textbf{2026-08-22}: Alp\"oge--Buckmaster's Lean verification (their statement); \textbf{2026-08-28}: OpenAI begins training the internal model used (the announcement); \textbf{2026-09-01}: OpenAI launches the agent effort after 'a rumor that two Millennium Prize problems had been resolved' (the announcement); \textbf{2026-09-03}: Tao's six-part post on the opportunity cost of an autonomous solution; Buckmaster writes to OpenAI (his statement); \textbf{2026-09-05}: OpenAI's agents 'arrive at their resolution', 88 hours in (the announcement); Tao's clarification: 'not aware of any significant developments'; \textbf{2026-09-06}: Lean formalization complete, 17 hours later (the announcement); the two calls with Buckmaster (his statement); \textbf{2026-09-08 03:58 UTC}: Alp\"oge--Buckmaster publish IPM, Boussinesq and Euler with smooth forcing, Lean at tristanbuckmaster/fluid\_lean (repository created 04:03 UTC); \textbf{2026-09-08 10:53 UTC}: openai/NavierStokesAndEuler created; its single commit at 10:57 UTC; the announcement the same day.}


\usepackage{array}
\newcolumntype{L}[1]{>{\raggedright\arraybackslash}p{#1}}
\newcommand{\lean}[1]{\texttt{#1}}
\newcommand{\Lp}[1]{L^{#1}}

\title{An audit of the 2026 Navier--Stokes finite-time blowup claim:\\
the formal statement read against Clay's wording and Mathlib's definitions,\\
the Lean proof built and its axioms asked of the kernel, and what remains open}
\author{\cmauthor}
\date{October 2026 --- draft v0.1 for review; every number interpolated from the record; machine-derived, not peer-reviewed; author list to be agreed}

\begin{document}
\maketitle

\begin{abstract}
On 8 September 2026 OpenAI published a \PaperPages-page proof, with a Lean~4 formalisation,
that for every positive viscosity a smooth compactly supported force drives the
three-dimensional incompressible Navier--Stokes flow from rest to unbounded velocity in
finite time, and stated that this establishes alternatives (C) and (D) of the Clay problem
statement \cite{OAIann,OAINS,Fef00}. This note records an audit made the following day from
a corpus pinned by digest: the formal statement read hypothesis by hypothesis against
Fefferman's wording and Mathlib's definitions; the \LeanLines-line repository built on a
laptop and the kernel asked for the axioms of each main declaration; Comparator run with a
second kernel; the paper mapped onto the certificate clause by clause; the writeup dissected
by \Readers{} independent readings and set against every classical obstruction; every
printed identity a computer can decide re-run with red controls. What
holds: the statement is Clay's with nothing added or dropped, the build passes on the three
standard axioms, both kernels accept, and no error or counterexample was found. Not decided:
the paper's finite-energy clause has no formal counterpart; the smallness parameter is
bounded numerically nowhere; \EulerChecksUnrun{} checks of the companion Euler paper are
listed and unrun. The register row reads \textsc{\LedgerVerdict}. Nothing is claimed about
the mathematics beyond what the kernel and the readings establish, nor about priority.
\end{abstract}

\section{Introduction}

The claim is Theorem~1.1 of \cite{OAINS}: for every $\nu>0$ there are a force
$f\in C_c^\infty(\R^3\times(0,\infty);\R^3)$, a compact $K\subset\R^3$ and smooth $(u,p)$
on $\R^3\times[0,1)$, supported in $K$, solving the Navier--Stokes equations from rest, with
$\sup_{0\le t<1}\norm{u(t)}_{\Lp2}<\infty$ and $\limsup_{t\uparrow1}\norm{u(t)}_{\Lp\infty}
=\infty$; consequently no smooth solution on $\R^3\times[0,\infty)$ with the same force and
datum has uniformly bounded kinetic energy. The paper says this is alternative (C) of the
Clay problem statement \cite{Fef00} and, through its Corollary~10.6, alternative (D). A
Lean~4 repository \cite{OAIrepo} formalises both; the announcement \cite{OAIann} says the
proof was produced by an internal model and formalised by another.

A formal certificate changes what an audit of such a claim can be: three questions become
decidable by machine and one by reading. Is the formal statement the theorem the prose
names; does the proof term check, and on which axioms; does the paper say what the
certificate says; and where the paper asserts rather than displays, does the certificate
discharge the assertion? The position throughout is the one the records state: \emph{the
Lean certificate is the claim; the paper is its human-readable evidence.} A gap in the paper
is a gap in the writeup unless the Lean shows the same, and a printed identity that failed
would refute the formula, not the theorem \cite{Records}. The contribution is not a theorem
but the grammar of the audit and its public record; every number below is interpolated from
a record by a tool that refuses when a record no longer supports a sentence.
Sections~\ref{sec:claim}--\ref{sec:dates} give the claim and the corpus, the audit, what the
record decides, what is not claimed, and the dates as the parties state them, weighed by
nobody here.

\section{What was claimed, and what was pinned}\label{sec:claim}

\paragraph{Clay's alternatives.} Fefferman's statement \cite{Fef00} asks for a proof of one
of four assertions; in its wording, (C) \emph{Breakdown of Navier--Stokes solutions on
$\R^3$}: take $\nu>0$ and $n=3$; then there exist a smooth, divergence-free vector field
$u^\circ(x)$ on $\R^3$ and a smooth $f(x,t)$ on $\R^3\times[0,\infty)$, satisfying (4), (5),
for which there exist no solutions $(p,u)$ of (1), (2), (3), (6), (7) on
$\R^3\times[0,\infty)$; and (D) the same on $\R^3/\Z^3$ with (8), (9) in place of (4), (5)
and (10), (11) in place of (6), (7), the errata adding that the pressure is periodic. The
numbered conditions are those of Table~\ref{tab:fidelity}.

\paragraph{The mechanism.} As the records describe it \cite{Records}: a self-similar
axisymmetric core with radial scale $\tau^{1/2}$, axial scale $\tau^{1/2-h}$ and azimuthal
and axial speeds $\tau^{-1/2-h}$, $\tau=1-t$, $0<h<1/100$, whose momentum residual is
unbounded in an annulus; two families of shear-amplified oscillatory pulses on an auxiliary
torus that supply the missing stress through their averaged quadratic flux; corrections at
every order in $q^{2h}$; an exact radial-heat exterior; cutoffs. The force is
\emph{defined} as the residual of the constructed flow and shown to extend smoothly through
$t=1$; uniqueness among smooth finite-energy solutions with a merely smooth pressure
transfers the growth to any competitor, in the forced-blowup program of C\'ordoba and
Mart\'{\i}nez-Zoroa \cite{CMZ}.

\paragraph{The corpus.} Everything read was pinned on \CorpusFetched{} with its digest:
\CorpusFiles{} files: the two papers, the Clay statement, the Lean bridge files, the Formal
Conjectures statement, the announcement and the concurrent papers and statements \cite{OAINS,OAIE,OAIann,OAIrepo,FC,AB}. The
repository \cite{OAIrepo} has one commit, \lean{\LeanCommit}, at \LeanCommitDate{}
\LeanCommitTime; it pins Lean \LeanToolchain{} \cite{LeanMathlib} and Mathlib at \lean{\MathlibRev}
\cite{LeanMathlib}; it holds \LeanFiles{} files and \LeanLines{} lines, of which \NSLines{} lines
in \NSFiles{} files are the Navier--Stokes library, \EulerLines{} in \EulerFiles{} the Euler
library and \ChalLines{} in \ChalFiles{} the challenge statements. Its
\lean{formalization.yaml} declares the review \emph{\FormalReview} and the automation
\emph{\FormalAutomation}, with no \lean{sorry} \cite{Records}.

\section{The audit}\label{sec:audit}

\subsection{The statement, hypothesis by hypothesis}\label{sec:statement}

The theorem the certificate proves was not written by the prover. The Comparator challenge
module is the Formal Conjectures statement of the Clay alternatives \cite{FC}, added by
DeepMind on \FormalConjAdded{} (pull request \#\FormalConjPR), last touched on
\FormalConjTouched{} at revision \lean{\FormalConjRev}, the revision OpenAI pinned and the
revision on \lean{main} on the day of the audit. OpenAI's copy differs from upstream in
\UpstreamDiffLines{} lines, every one an \lean{open} or a notation, and the solution side
re-declares the same definitions under the same names, byte-identical once comments are
stripped \cite{Records}. Comparator compares every constant reachable from the theorem's
type for exact equality, so the only question is whether that statement is Clay's.

Both theorems have the shape $\exists\,u_0\,f,\ H(u_0)\wedge H(f)\wedge\neg\exists\,v\,p,\
S(v,p)$: the prover must \emph{establish} $H$ for the constructed data and \emph{exclude}
every $(v,p)$ in $S$. A junk value or a weak definition inside $S$ only enlarges the class
and strengthens the theorem; a gap could hide only in an $S$ stronger than Clay's (6)(7) or
(10)(11), excluding competitors Clay admits, or in an $H$ weaker than Clay's (4)(5) or
(8)(9), admitting data Clay rejects. Every row of Table~\ref{tab:fidelity} is read in that
direction, with the Mathlib definition located at the pinned revision \cite{Records}.

\begin{table}[t]\centering\footnotesize
\begin{tabular}{L{3.3cm}L{8.6cm}L{2.6cm}}\toprule
Clay & Lean (Formal Conjectures) and the Mathlib definition at \lean{\MathlibRev} & direction\\\midrule
(1) $\partial_tu+(u\cdot\nabla)u=\nu\Delta u-\nabla p+f$ on $\R^3\times[0,\infty)$ &
\lean{derivWithin (v x $\cdot$) (Ici 0) t + fderiv $\mathbb{R}$ (v $\cdot$ t) x (v x t) = nu $\bullet$ $\Delta$ (v $\cdot$ t) x - gradient (p $\cdot$ t) x + f x t}, $\forall x$, $\forall t\ge0$: the one-sided time derivative at $t=0$ within $[0,\infty)$; \lean{fderiv} applied to $v$ is $(v\cdot\nabla)v$; \lean{$\Delta$} the ordinary Laplacian (trace of the second derivative); \lean{gradient} the Riesz representative of \lean{fderiv} & same \\
(2) $\operatorname{div}u=0$ & \lean{$\nabla\!\cdot$ (v $\cdot$ t) x = 0}, the trace of the Jacobian & same \\
(3) $u(x,0)=u^\circ$ & \lean{v x 0 = u$_0$ x} & same \\
(6) $p,u\in C^\infty(\R^3\times[0,\infty))$ & \lean{ContDiffOn $\mathbb{R}$ $\infty$ ($\upharpoonleft$v) (univ $\times^{s}$ Ici 0)}, likewise $p$: $C^\infty$ up to $t=0$ in the Taylor-series sense, not analytic & same \\
(7) $\int|u|^2dx<C$ for all $t$ & \lean{integrable : $\forall$ t $\ge$ 0, MemLp ($\|$v $\cdot$ t$\|$) 2} with \lean{globally\_bounded\_energy : $\exists$ E, $\forall$ t $\ge$ 0, $\int$ $\|$v x t$\|$$^2$ < E}; the Bochner integral of a non-integrable function is $0$, which the first field forbids & same \\
(10) $u$ periodic; errata: $p$ periodic; (11) smooth & \lean{IsOnePeriodic (v $\cdot$ t)}, \lean{IsOnePeriodic (p $\cdot$ t)} for $t\ge0$, along \lean{EuclideanSpace.single i 1}; smoothness as (6) & same as Clay plus the errata \\
(4) $u^\circ$ smooth, divergence-free, $|\partial^\alpha u^\circ|\le C(1+|x|)^{-K}$ $\forall\alpha,K$ & \lean{InitialVelocityConditionDecay}: divergence-free, \lean{ContDiff $\mathbb{R}$ $\infty$ u$_0$}, $\forall m\,K\ \exists C$, \lean{$\|$iteratedFDeriv $\mathbb{R}$ m u$_0$ x$\|$ $\le$ C/(1+$\|$x$\|$)\^{}K}; the operator norm of the $m$-th derivative is comparable to the largest order-$m$ partial & equivalent; the paper takes $u^\circ=0$ \\
(5) $f$ smooth on $\R^3\times[0,\infty)$, $|\partial_x^\alpha\partial_t^mf|\le C(1+|x|+t)^{-K}$ & \lean{ForceConditionDecay}: \lean{ContDiffOn $\mathbb{R}$ $\infty$ ($\upharpoonleft$f) (univ $\times^{s}$ Ici 0)} and the same bound on \lean{iteratedFDerivWithin} of total space-time order $m$ in place of Clay's split $(\alpha,m)$ & equivalent; a $C_c^\infty$ force satisfies it trivially \\
(8) $u^\circ$, $f$ periodic in $x$; (9) $|\partial^\alpha\partial^mf|\le C(1+t)^{-K}$ & \lean{InitialVelocityConditionPeriodic}; \lean{ForceConditionPeriodic} with the time-decay bound on \lean{iteratedFDerivWithin} & same \\
\bottomrule\end{tabular}
\caption{Clay's hypotheses against the formal statement, one row per hypothesis
(\FidelityRows{} rows; \FidelitySame{} read \emph{same}, \FidelityEquivalent{}
\emph{equivalent}). The first six rows are the competitor class $S$, the last three the
data hypotheses $H$. Transcribed from the record \cite{Records}.}\label{tab:fidelity}
\end{table}

\begin{proposition}[statement fidelity; the record's verdict]\label{prop:fidelity}
The formal statements of (C) and (D), \lean{navier\_stokes\_breakdown\_R3} and
\lean{navier\_stokes\_breakdown\_periodic}, are Clay's with no hypothesis added to the
competitor class and none dropped from the data; in particular no decay, growth or
integrability condition is imposed on the pressure or on $\nabla v$ beyond smoothness, so the
formal uniqueness argument must handle every competitor Clay admits. The dimension is
\lean{Fin~3} and $\nu$ is any positive real.
\end{proposition}

A negation is cheap if the class it negates is empty for a trivial reason, so a witness was
written for this audit and compiled against OpenAI's challenge module: the fluid at rest
satisfies every field of both competitor structures and the zero data all four data
conditions (\lean{\NonVacuityFile}, \NonVacuityDecls{} declarations, \NonVacuityErrors{}
errors, the three standard axioms for \lean{classes\_inhabited}) \cite{Records}. It shows
the definitions are satisfiable, not that they admit every classical solution, which is
Table~\ref{tab:fidelity}'s job.

\subsection{The build, the axioms and the second kernel}\label{sec:build}

{\sloppy The repository was built on \BuildDate{} on an ordinary laptop (\BuildMachine, \BuildCpus{}
cores, \BuildMemory~GB, \BuildOS, \BuildArch) with the pinned toolchain: all \BuildJobs{}
modules in \BuildMinutes{} minutes after the Mathlib cache (\BuildJobsTotal{} Lake jobs in
all; two legs, \BuildLegOne{} modules at \BuildLegOneThreads{} threads then \BuildLegTwo{} at
\BuildLegTwoThreads, after eight parallel Lean processes drove the machine into swap), exit
\BuildExit, no errors \cite{Records}. A scratch file importing the two solution
modules then asked \lean{\#print axioms} of the \AxiomDecls{} declarations
\lean{formalization.yaml} names, and the kernel reported exactly \AxiomList{} for each. In
code lines, comments blanked, the proof libraries contain no \lean{sorry} beyond the
\HygSorryChal{} intentional placeholders of the challenge files, no axiom declaration, no
\lean{native\_decide}, \lean{unsafe}, \lean{opaque} or \lean{partial def}, no custom syntax
and no heartbeat override (\HygHeartbeatComments{} docstrings record earlier overrides
removed).\par}

Then the judge. Comparator \cite{Comparator}, the Lean FRO's tool for exactly this
situation, exports both sides with \lean{lean4export}, compares every constant reachable
from each theorem's type for exact equality, walks the axioms, and replays the proof through
Lean's kernel and through nanoda, a kernel written independently of Lean's in Rust and built
from source for this run. On the \ComparatorConfigs{} configurations, Navier--Stokes and
Euler, all \ComparatorKernelAccepts{} kernel replays accept and \ComparatorAccepted{} of
\ComparatorConfigs{} return \emph{``Your solution is okay!''} with exit $0$, in
\ComparatorMinutesNS{} and \ComparatorMinutesEuler{} minutes. One component did not run and
the record says so: Comparator's Landlock process isolation is Linux-only and a pass-through
stand-in took its place; the statement comparison, the axiom walk and both kernel replays
were performed, the isolation was not, and on a machine where both sides come from one
pinned commit it guards against nothing this audit needed. What the build establishes is
bounded by its trusted base, the two kernels, Mathlib's definitions at the pinned revision,
the elaborator and the hardware: the proof terms exist and check against the third party's
statement, and nothing about the paper.

\subsection{The paper against the certificate}\label{sec:map}

Each clause of the stated theorem was located in the formal library or named as absent
(Table~\ref{tab:map}; \PaperVsLeanRows{} rows in the record, \PaperClausesNotFormal{} of them
not formal) \cite{Records}.

\begin{table}[t]\centering\footnotesize
\begin{tabular}{L{2.5cm}L{4.0cm}L{5.6cm}L{2.4cm}}\toprule
item & the paper & the Lean & status\\\midrule
the two theorems & Theorem~1.1 (C), Corollary~10.6 (D) & \lean{navier\_stokes\_breakdown\_R3}, \lean{\_periodic}, through \lean{ComparatorBridge} adapters; solution-side definitions byte-identical to the challenge & proved, subject to the build record \\
uniqueness, Lemma~10.5 & $v$ smooth, $\Lp\infty_t\Lp2_x$, zero datum, same force $\Rightarrow v=u$; no growth condition on the pressure & \lean{classical\_uniqueness\_on\_Icc} (\lean{R3/WholeSpaceUniqueness.lean:30}): competitor smooth, divergence-free, uniformly finite energy, pressure only smooth; pressure recovered (\lean{R3/PressureRecovery.lean:419}), flux bounded uniformly in the radius (\lean{R3/PressureFlux.lean:576}) & covers every competitor Clay admits \\
the viscosity rescaling & (10.22): space rescaled, singular time $1$ & time rescaled, $f_\nu(x,t)=\nu^2f(\nu t,x)$ (\lean{ComparatorR3Theorem.lean:19}), singular time $1/\nu$ & equivalent; the statement names no time \\
the uniform energy bound & Theorem~1.1's $\sup_{t<1}\norm{u(t)}_{\Lp2}<\infty$; Lemma~10.4; the announcement's ``its energy remains finite through the entire dynamics'' & \lean{uniform\_finite\_energy} (\lean{R3/CompactEnergy.lean:343}) proves it for any smooth compactly supported force and is used by nothing; the structure carrying \lean{energy\_bounded} has no producer; the candidate's \lean{Properties} structure has ten fields and no energy field & \textsc{not formal}; every other clause of Theorem~1.1 is \\
the force's support in time & $f\in C_c^\infty(\R^3\times(0,\infty))$: zero near $t=0$ & smooth for $t\ge0$, compact spatial support, zero for $t\ge T$; no field states $f=0$ near $t=0$ & weaker in Lean; the challenge's force condition does not require it \\
the growth & (10.20)--(10.21): $u_\theta=\tau^{-A}(e_0+O(\tau^{2h}))$ along a named off-axis path & \lean{SpeedUnboundedAtOne}: $\forall M\,\delta\ \exists t\,x$, $1-\delta<t<1$, $M<\norm{u(t,x)}$; on the torus also $\norm{u(t,0)}\to\infty$ on the axis & pointwise $\limsup\norm u_\infty=\infty$; no rate \\
the numbering & Propositions 5.5, 7.5, 9.6, 9.9, Theorem~3.1 & docstrings cite a manuscript numbering absent from the published paper & the map is by content \\
\bottomrule\end{tabular}
\caption{The paper mapped onto the certificate, from the record \cite{Records}.}\label{tab:map}
\end{table}

One difference bears on the announcement, which says the fluid's ``energy remains finite
through the entire dynamics, from rest to the formation of the singularity'' \cite{OAIann}.
Theorem~1.1 states it and Lemma~10.4 proves it in a few lines, re-derived in the reading;
the repository contains the general lemma, nothing invokes it, and
\lean{breakdownStatement}, the file's own ``full assertion of Theorem~1.1'', is never
proved. Clay's (C) does not need the clause, since a global smooth competitor is excluded by
uniqueness on every $[0,T]$, $T<1$, where the candidate's energy is finite by compact support
alone; so the formal result stands as Clay's and the physically striking clause rests on the
paper. The Euler theorem in the same repository carries its energy bound
(Section~\ref{sec:euler}); the record reads the asymmetry as an oversight rather than an
obstacle, and this note repeats that reading as a reading.

\subsection{The writeup, read}\label{sec:read}

The paper was taken apart the day after it appeared by \Readers{} independent readings,
\SectionMemos{} section memos (\S3 with \S10, \S4 with Appendix~B, \S5 with Appendix~A,
\S6 with \S7, \S8 and \S9 with Appendix~C) and one map of the Lean library against all of
them, each from the pinned PDF, written in the machine-assisted setting the author's
disclosure describes and consolidated into one record \cite{Records}. Each reading
transcribed the statements, traced the order in which every parameter is fixed and
re-derived what it could; no circularity and no error were found. What was verified by hand holds as printed in all \VerifiedRows{} groups of
the record: the closing argument of \S10 term by term; the leading equations of \S4
re-derived in the similarity coordinates; the heat exterior of Appendix~A with Lemma~A.8's
two zero-integral identities; the identities of \S\S6--7; the moment system of \S8 and the
loop of Lemma~C.1 on \ConeLoopTrials{} random relaxed-cone data, where choosing $\delta_L$
from the base point alone breaks the cone in about \ConeBreakPct\% of trials, so the
ordering ``$\delta_L$ after $\mu_{\max}$, uniformly'' is load-bearing; and the coefficient
equations of \S5, by computer (Section~\ref{sec:probes}).

What the paper asserts rather than displays was listed and ranked by how much rests on it,
\WeakPoints{} points in all \cite{Records}. Point~0 is the subject of
Section~\ref{sec:obstructions}; the other \WeakPointsAfterZero, in the record's order, are
the endpoint regularity of the local field up to $\tau=0$ in Proposition~9.9 (one paragraph,
no estimate displayed, on which the force's smoothness at $t=1$ rests);
the uniformity in $\eta$ of the flatness (3.4) that (10.9) needs; the admissibility of
next-cycle sources in Proposition~9.3(i), the one unstated induction hypothesis of \S9; the
smallness parameters, every one ``sufficiently small'' and none numeric, the printed
$h<1/100$ not being the operative bound (Lemma~4.8 needs $h<\min\{\lambda,e^{-T_d}\}$ with
$T_d=e^{M_d}+10$, the derivative scale of \S6 exceeds one only from dyadic level about
$3\times10^{8}$, and flatness to order $N$ needs more than $1000(N+K_m)$ cycles and about
$10^{7}N$ integrations by parts); the Lipschitz bounds of Proposition~B.2; the inequality
$2+h<a$ of Lemma~4.9 on the terminal collar; the continuation of Proposition~B.5; the
analyticity region of Lemma~5.1 against Theorem~4.6(i); the coefficient equations
(5.3)--(5.6), since discharged by computer algebra; the two correction routes of \S7 and
\S8.5; the hypothesis $|T_0|\ge c\zeta$ of Proposition~C.3, never cited in \S9; and the
overloading of $q_*$. None is an error. \Slips{} bookkeeping slips were found, each
harmless with a positive margin (a dropped $-\kappa_s$ in the zero harmonic of two-wave
products has margins $0.08$ and $0.4$ against $\kappa_s=10^{-5}$). One is a warning to
anyone auditing from extracted text: the cutoff family of Lemma~10.5 on p.~122 extracts as
the dilations $\varphi_{4R},\varphi_{8R}$, under which three inequalities are false, while
the glyphs are the powers $\varphi_R^4,\varphi_R^8$, under which every step is exact; no
inequality is a finding until its glyphs have been seen.

\subsection{The asserted steps, traced into the kernel}\label{sec:traced}

Each asserted step was traced into the Lean library at the pinned commit and asked a harsher
question than ``is it discharged'': whether the formal statement there is \emph{weaker}
than the paper's claim \cite{Records}. \TracedFull{} are discharged by a proved theorem at
full strength, \TracedStronger{} of them stronger than the paper: the endpoint regularity is
a genuine $C^\infty$ extension across $t=1$ (\lean{OneSidedExtension},
\lean{JointResidualLimits.lean:72}) and the flatness a joint space-time limit,
$\eta$-uniform by construction (\lean{VanishingJointJets},
\lean{JointResidualLimits.lean:84}), exactly the uniformity the reading found missing.
\TracedAltered{} are discharged with something altered, among them the smallness
parameters: $h$ is not a number in Lean but \lean{Classical.choice} of a nonempty profile
datum, constrained by $2h<\lambda<1/10$ and the hypothesis $h\le1/1000$
(\lean{NaturalAxisData.lean:41}); $N$, $\Lambda$, $C$, $q_*$ and the cycle schedule are
existentials; the only numeric constants are $\kappa=10^{-5}$,
$\Lambda_{\mathrm{chart}}=4-\sqrt2$, $X_{\mathrm{big}}=100$ and $X_i=110$; the Lipschitz
machinery of Proposition~B.2 and the analyticity region are proved as machinery, their
instantiation for the actual profile not traced to a line. \TracedNone{} point, Lemma~4.9's
collar inequality, has no counterpart, the formal tail being a different plateau
construction; \SlipsProseOnly{} of the slips are prose only, and two of the paper's
arithmetic slips are absent from the formal object.

Two more answers were read from the source. The construction's non-axisymmetry, which
Section~\ref{sec:obstructions} needs, is a theorem, \lean{angularMode\_ne\_zero}
(\lean{ActualInitialization.lean:52}), not an assumption. And the growth witness,
\lean{FinalSlowBase.axis\_tendsto}, is the \emph{axial} component on the axis along
$t\mapsto(t,0)$, not the paper's off-axis azimuthal path (10.20)--(10.21); both give
$\limsup\norm u_\infty=\infty$, and no Lean theorem bounds the swirl $r\,u_\theta$ from below.

\subsection{The adversarial pass}\label{sec:obstructions}

The question an audit of a singularity must ask is whether a theorem says this cannot
happen. \ObstructionRows{} classical obstructions were set against the construction's own
exponents \cite{Records}. Of these,
\ObstructionsMet{} are necessary conditions and are met: Serrin's
$\int\norm u_\infty^2dt=\infty$ \cite{Serrin}, with $\norm u_\infty^2\asymp\tau^{-1-2h}$;
Escauriaza--Seregin--\v{S}ver\'ak's $\norm{u(t)}_{\Lp3}\to\infty$ \cite{ESS}, met slowly, as
$\tau^{-4h/3}$; the Beale--Kato--Majda integral \cite{BKM}; Leray's lower bounds
\cite{Leray}; finite energy and finite dissipation, the dissipation rate $\tau^{-1/2-3h}$
being integrable exactly when $h<1/6$; and the Caffarelli--Kohn--Nirenberg measure bound
\cite{CKN}, one point at one time. The other \ObstructionsEvaded{} are exclusions and are
evaded: the type-I exclusions for axisymmetric flows \cite{CSYT,KNSS} twice, the flow being
type II by the factor $\tau^{-h}$ and not axisymmetric, and the classical no-swirl
regularity, the force injecting swirl.

The third is load-bearing and unstated. For an axisymmetric flow the swirl $\Gamma=r\,u_\theta$
satisfies
\[
\partial_t\Gamma+u_r\partial_r\Gamma+u_z\partial_z\Gamma
=\Gamma_{rr}-\Gamma_r/r+\Gamma_{zz}+r f_\theta ,
\]
a drift--diffusion equation with no zeroth-order term, derived in the battery from the
$\theta$-momentum equation and checked symbolically; so from rest with a bounded compactly
supported force, $\sup|\Gamma(t)|\le\int_0^t\sup|rf_\theta|$ and $|u_\theta|\le C/r$ for every
$t<1$. The paper's leading field has $r\,u_\theta^{(0)}=q^{-h}H$, in its own words on p.~27,
unbounded as $q\downarrow0$. The axisymmetric background alone is therefore impossible as a
forced flow from rest; the theorem survives because the physical field is not axisymmetric,
the pulses carrying nonzero integer angular frequencies (p.~12) whose Reynolds flux breaks
the maximum principle for the angular mean. The paper does not say this anywhere; it is the
sentence a referee would ask for, and it is why $h>0$ is forced rather than chosen. The
formal object evades the obstruction by \lean{angularMode\_ne\_zero}
(Section~\ref{sec:traced}). This is a reading, not a refutation: the record is explicit that
no obstruction refutes the construction.

\subsection{The printed identities, re-run}\label{sec:probes}

The proof carries no numerics and fixes no constant to a number, so there is nothing for an
interval certifier to certify. What the paper prints are identities, and each was re-run
from the formula as printed by \ProbeScripts{} probe scripts carrying \ProbeReds{} red
controls, a deliberately wrong variant of each test that must be rejected, \ProbeRedsDead{}
of which failed to fire \cite{Records}. Symbolically: the similarity-coordinate derivatives
of Lemma~4.1, the commutator kernel norm $16\pi/R$ of Lemma~10.5, the derivative count
(10.12), the rescalings (10.22)--(10.23) and of Corollary~10.6, the exterior
$q$-cancellation, the cutoff bound (10.3), and the exact exterior swirl
$K=(r^2/2)^{-A}H(4\tau/r^2)$ against its equation (A.37). At \HeatDigits{} digits: the
integral $H$ of (A.32) against (A.37), worst relative residual \HeatResidual, and against
its Taylor coefficients (A.35) to order \HeatTaylorOrder. By exact algebra: the expansion
(5.1) substituted into the axisymmetric operator and compared term by term with the printed
coefficient equations (5.2)--(5.6), the sparsity and block nilpotency of $A_1$ in (5.7) and
the divisibility of $\Omega_k$ by $X$, \CoeffIdentities{} identities. In exact rationals:
the admissible stress cone of \S4.3, where Lemma~4.5's square-root-free form (4.22) agrees
with the relaxed condition on a sweep of \ConeSweepPoints{} points (\ConeAdmissible{}
admissible, \ConeNotAdmissible{} not), a planted sign flip disagrees at
\ConeSignFlipDisagree{} of them, and Proposition~7.5's printed amplitudes are real exactly on
the reference cone at \ConeRefPoints{} points (\ConeRefInside{} inside); and the swirl
equation above. In floating point, deciding nothing, as the record says: the axis
initial-value problem of Appendix~B at representative data the paper never fixes, whose
solved profile tends to the closed form of Proposition~B.2 as $1/\Lambda$, with (B.17)
holding from $\Lambda=\AxisLambdaFrom$ on. The record also keeps this battery's own first
bug, a wrong derivative factor ($(-h)_m$ for $(-1)^m(h)_m$) that made a first run report $H$
failing its equation; it is now a red control, so that a scratch log reading \textsc{fail}
against a Millennium claim is not mistaken for a finding.

\subsection{The companion theorem}\label{sec:euler}

The same repository formalises Theorem~1.1 of \cite{OAIE}: there is
$u_0\in C^\infty_{c,\sigma}(\R^3)$ with $0<T^*(u_0)<\infty$ whose smooth unforced Euler
solution has $\limsup_{t\uparrow T^*}\norm{\nabla u(t)}_{\Lp\infty}=\infty$ and
$\int_0^{T^*}\norm{\operatorname{curl}u(t)}_{\Lp\infty}dt=\infty$, a breakdown with no
force, read in full \cite{Records}. {\sloppy Of its two formal shapes,
\lean{exists\_compact\_smooth\_euler\_singularity} asserts a nonzero smooth compactly
supported datum, $0<T^*\le1$, a uniform energy bound on the whole lifespan, maximality as
an equivalence, infinite $\limsup$ of the $C^1$ norm and infinite Beale--Kato--Majda
integral at $T^*$, and no global smooth finite-energy solution; \lean{euler\_breakdown\_R3}
is the strictly weaker non-existence form, equivalent to blowup for these data only because
the formal class imposes square integrability and bounded energy.\par}

The step where such arguments die is the
transfer from the approximants $U_j$ to the limiting datum, Euler's stability constant
$\exp(C\int\norm{\nabla U_j}_\infty)$ being the diverging quantity; the reading finds \S6.2
complete as written, because it compares each $U_j$ with the hypothetical smooth solution
of the limiting datum under the contradiction hypothesis $T^*(u_0)>T_\infty$, so the
Gronwall constant is $C\sup\norm u_{H^4}$, fixed before $j$ is chosen. The reading's
sharpest candidate error, (3.16) on p.~13 apparently printing $(k\cdot C_*)^{n+1}$ with
$C_*=80$, false at $n=0$, was \textsc{withdrawn}: the page rendered at 150~dpi shows
$(k^{C_*})^{n+1}(n!)^2$; the extraction had flattened the superscript, the second such
artifact of this audit. The record lists \EulerChecksUnrun{}
checks the reading transcribed and did not run: the frozen-frame WKB model as an exact ODE
system, every scale inequality of \S5.5 and (6.1)--(6.2) in exact logarithmic arithmetic,
and the symbolic identities of \S3 including the trace check that fixes the $-2$ in (3.14).
The Euler result is audited separately and not folded into the Navier--Stokes verdict.

\section{What the record decides}\label{sec:decides}

The register row \lean{\LedgerId}, recorded \LedgerRecorded, reads \textsc{\LedgerVerdict}:
the Lean certificate holds for Clay's (C) and (D), and the finite-energy clause the paper's
Theorem~1.1 states is not in the formal statement \cite{Records}. Exactly:

\begin{enumerate}
\item \emph{Held.} The formal statements are Clay's (C) and (D) with nothing added to the
competitor class and nothing dropped from the data, written by a third party before the
proof existed; the competitor classes are inhabited (Section~\ref{sec:statement}).
\item \emph{Held.} All \BuildJobs{} modules build on this machine with the pinned toolchain;
the kernel reports exactly \AxiomList{} for each of the \AxiomDecls{} main declarations;
Comparator accepts both configurations through Lean's kernel and nanoda
(Section~\ref{sec:build}).
\item \emph{Held.} Every clause of Theorem~1.1 except the uniform energy bound and the
vanishing of the force near $t=0$ has a formal counterpart, and the formal uniqueness covers
every competitor Clay admits (Section~\ref{sec:map}).
\item \emph{Held, as readings.} No error, circularity or counterexample was found in
\Readers{} independent readings; \TracedFull{} of the \WeakPointsAfterZero{} asserted steps
are discharged in Lean at full strength, \TracedStronger{} more strongly than the paper
states; no classical obstruction refutes the construction; all \ProbeScripts{} probe scripts
pass with all \ProbeReds{} red controls firing; the Euler transfer step reads complete
(Sections~\ref{sec:read}--\ref{sec:euler}).
\item \emph{Not decided.} The paper's $\sup_{t<1}\norm{u(t)}_{\Lp2}<\infty$ and the force's
vanishing near $t=0$ are paper claims with no formal counterpart; the first rests on the
paper's Lemma~10.4, re-derived by hand.
\item \emph{Not decided.} The smallness parameter $h$ is bounded numerically in neither
artifact; the proof is asymptotic in a parameter no one has quantified, and nothing here
makes it effective.
\item \emph{Not decided.} The instantiation of Proposition~B.2's machinery for the actual
profile was not traced to a line; Lemma~4.9's collar inequality has no formal counterpart;
the paper names no sentence for the swirl obstruction it evades; the \EulerChecksUnrun{}
Euler checks are listed and unrun; Comparator's process isolation did not run.
\item \emph{Not decided.} The mathematics of the writeup beyond what the kernel and the
readings establish. An audit is not a referee report, and the readings were not peer review.
\end{enumerate}

\section{What is not claimed}\label{sec:not}

Nothing is claimed about the correctness of the mathematics beyond items 1--4 above: the
kernel establishes that the proof terms check against the third party's statement, the
readings that nobody here found an error, and the verdict on the theorem is the build's and
the kernels', not the reading's. No numerical computation here certifies anything: the
paper contains no numerics and none of its constants is a number. Nothing is claimed about
priority: the record holds the dates each party states (Section~\ref{sec:dates}) and
adjudicates nothing, and the concurrent results \cite{AB} are different theorems from the
one audited here. Nothing is claimed about the prize: the papers are published and not
peer-reviewed, the Clay Mathematics Institute had said nothing at the time of the record,
and OpenAI has said it does not intend to claim it \cite{OAIann}. The word \emph{certified}
is used here only for what a kernel checked.

\section{The dates, as the parties state them}\label{sec:dates}

The record holds \DateRowsCount{} dates, each from the document named in it: OpenAI's
announcement \cite{OAIann}, Buckmaster's statement and the concurrent papers \cite{AB},
Tao's posts \cite{Tao}, or a repository's own timestamps \cite{OAIrepo,FC}. Nothing here
weighs them \cite{Records}.

{\footnotesize\DateList\par}

\cmrepro{\sloppy The records are in \texttt{corpus/navier-stokes/} \cite{Records}:
\texttt{lean-repo.json} is computed from a clone at \texttt{\LeanCommit} by
\texttt{tools/pin-navierstokes-lean.js}, \texttt{build.json} by
\texttt{tools/record-navierstokes-build.js}, and \texttt{probes.json} by
\texttt{python3 instruments/navierstokes/battery.py}, which re-runs the \ProbeScripts{}
probe scripts with their \ProbeReds{} red controls. The macros of this document are
written by \texttt{node tools/paper-numbers/navier-stokes-audit.js} from those records at
repository commit \texttt{\RepoCommit}; the timing fields of \texttt{probes.json}, which
change on every run, are never read. Rebuilding the Lean repository needs the clone and
toolchain the two tools name and about \BuildMinutes{} minutes on an eight-core machine;
nothing here depends on re-running it, and every number above is read from the record of
the run of \BuildDate.}

\begin{thebibliography}{99}\footnotesize
\bibitem{Fef00} C.~L. Fefferman, Existence and smoothness of the Navier--Stokes equation,
Clay Mathematics Institute, 2000, with errata,
\url{https://www.claymath.org/wp-content/uploads/2022/06/navierstokes.pdf}, pinned \CorpusFetched.
\bibitem{OAINS} OpenAI, Finite time blowup for Navier--Stokes, \PaperPages~pp., 8 September
2026, \url{https://cdn.openai.com/pdf/32d9f210-8b73-45e0-91bc-82a30aef8a9a/navier-stokes.pdf}, pinned sha256 \texttt{\PaperSha}\ldots
\bibitem{OAIE} OpenAI, Finite time blowup for the Euler equation, \EulerPages~pp., 8 September
2026, \url{https://cdn.openai.com/pdf/315b36cd-ec98-4023-8342-93345194ece1/euler.pdf}.
\bibitem{OAIann} OpenAI, On the Navier--Stokes Millennium Prize Problem, 8 September 2026,
\url{https://openai.com/index/navier-stokes-solution/}, pinned \CorpusFetched{} via the Internet Archive.
\bibitem{OAIrepo} OpenAI, NavierStokesAndEuler: Lean 4 formalizations,
\url{https://github.com/openai/NavierStokesAndEuler}, commit \texttt{\LeanCommit}, \LeanCommitDate.
\bibitem{FC} Google DeepMind, Formal Conjectures,
\texttt{FormalConjectures/Millenium/NavierStokes.lean},
\url{https://github.com/google-deepmind/formal-conjectures}, revision
\texttt{\FormalConjRev}; the statement added \FormalConjAdded{} in \#\FormalConjPR.
\bibitem{LeanMathlib} L. de Moura, S. Ullrich, The Lean 4 theorem prover and programming
language, in: Automated Deduction -- CADE 28, LNCS 12699, Springer, 2021, 625--635; The
mathlib Community, The Lean mathematical library, in: Proc.\ CPP 2020, ACM, 367--381;
Mathlib revision \texttt{\MathlibRev} as pinned.
\bibitem{Comparator} Lean FRO, Comparator, \url{https://github.com/leanprover/comparator},
revision \texttt{\ComparatorRev}; lean4export,
\url{https://github.com/leanprover/lean4export}, revision \texttt{\LeanExportRev}; nanoda
(\texttt{ammkrn/nanoda\_lib}), a Lean kernel written in Rust, built from source for the run.
\bibitem{CMZ} D. C\'ordoba, L. Mart\'{\i}nez-Zoroa, Finite time singularities of smooth
solutions for the 2D incompressible porous media (IPM) equation with a smooth source,
arXiv:2410.22920v3 (2025).
\bibitem{AB} L. Alp\"oge, T. Buckmaster, Blowup for the Boussinesq equations with smooth
forcing, and Blowup for the Euler equations with smooth forcing; with M.~P. Coiculescu,
Extending the C\'ordoba--Mart\'{\i}nez-Zoroa IPM blow-up to uniformly space-time smooth
forcing; T. Buckmaster, statement of 8 September 2026; all at
\url{https://cims.nyu.edu/~tristanb/}, pinned \CorpusFetched.
\bibitem{Tao} T. Tao, posts of 3 and 5 September 2026,
\url{https://mathstodon.xyz/@tao/117207849921390904}, pinned \CorpusFetched.
\bibitem{Leray} J. Leray, Sur le mouvement d'un liquide visqueux emplissant l'espace, Acta
Math.\ 63 (1934), 193--248.
\bibitem{Serrin} J. Serrin, On the interior regularity of weak solutions of the Navier--Stokes
equations, Arch.\ Rational Mech.\ Anal.\ 9 (1962), 187--195.
\bibitem{CKN} L. Caffarelli, R. Kohn, L. Nirenberg, Partial regularity of suitable weak
solutions of the Navier--Stokes equations, Comm.\ Pure Appl.\ Math.\ 35 (1982), 771--831.
\bibitem{BKM} J.~T. Beale, T. Kato, A. Majda, Remarks on the breakdown of smooth solutions
for the 3-D Euler equations, Comm.\ Math.\ Phys.\ 94 (1984), 61--66.
\bibitem{ESS} L. Escauriaza, G. Seregin, V. \v{S}ver\'ak, $L_{3,\infty}$-solutions of
Navier--Stokes equations and backward uniqueness, Russian Math.\ Surveys 58 (2003), 211--250.
\bibitem{KNSS} G. Koch, N. Nadirashvili, G. Seregin, V. \v{S}ver\'ak, Liouville theorems for
the Navier--Stokes equations and applications, Acta Math.\ 203 (2009), 83--105.
\bibitem{CSYT} C.-C. Chen, R.~M. Strain, H.-T. Yau, T.-P. Tsai, Lower bound on the blow-up
rate of the axisymmetric Navier--Stokes equations, Int.\ Math.\ Res.\ Not.\ IMRN 2008, Art.\
ID rnn016.
\bibitem{Records} cert-machine, the records of the audit in \texttt{corpus/navier-stokes/}
(\texttt{MANIFEST.json}, \texttt{lean-repo.json}, \texttt{build.json},
\texttt{nonvacuity.json}, \texttt{fidelity-clay-vs-lean.md}, \texttt{dissection.md},
\texttt{audit.json}, \texttt{euler.md}, \texttt{probes.json}), all of 2026-09-09; the
register row \texttt{\LedgerId} in \texttt{certs/claims-ledger.json}; and the page built
from them, \url{https://carlostoledo.co/reports/navier-stokes.html}.
\end{thebibliography}

\end{document}
